% Part: turing-machines
% Chapter: machines-computations
% Section: variants

\documentclass[../../../include/open-logic-section]{subfiles}

\begin{document}

\olfileid[pt]{tur}{mac}{var}
\olsection{Variantes de Máquinas de Turing}

Na verdade, há muitas maneiras possíveis de definir máquinas de Turing, das
quais a nossa é apenas uma.  De certas formas, nossa definição é mais liberal
do que outras. Permitimos alfabetos finitos arbitrários; uma definição mais
restrita poderia permitir apenas dois símbolos de fita, $\TMstroke$
e~$\TMblank$.  Permitimos que a máquina escreva um símbolo na fita
e se mova ao mesmo tempo; outras definições permitem escrever ou
mover.  Permitimos a possibilidade de escrever sem mover a cabeça da fita;
outras definições omitem a ``instrução'' $\TMstay$.  Em outros aspectos,
nossa definição é mais restritiva. Assumimos que a fita é infinita em apenas
uma direção; outras definições permitem que a fita seja infinita tanto para a
esquerda quanto para a direita. Na verdade, pode-se até mesmo permitir
qualquer número de fitas separadas, ou até mesmo uma grade infinita de
quadrados.  Representamos o conjunto de instruções da máquina de Turing por
uma função de transição; outras definições usam uma relação de transição em
que a máquina tem mais de uma instrução possível em qualquer situação dada.

Essa última flexibilização da definição é particularmente interessante.
Em nossa definição, quando a máquina está no estado~$q$ lendo o
símbolo~$\sigma$, $\delta(q, \sigma)$ determina qual é o novo símbolo, o
novo estado e a nova posição da cabeça da fita.  Mas se permitirmos que o
conjunto de instruções seja uma relação entre pares estado-símbolo atuais
$\tuple{q, \sigma}$ e novas triplas estado-símbolo-direção
$\tuple{q', \sigma', D}$, a ação da máquina de Turing pode não ser
exclusivamente determinada --- a relação de instrução pode conter tanto
$\tuple{q, \sigma, q', \sigma', D}$ quanto
$\tuple{q, \sigma, q'', \sigma'', D'}$.
Nesse caso, temos uma máquina de Turing \emph{não determinística}.  Estas
desempenham um papel importante na teoria da complexidade computacional.

Há também diferentes convenções para quando uma máquina de Turing para:
dizemos que ela para quando a função de transição é indefinida; outras
definições exigem que a máquina esteja em um estado de parada especial
designado. Explicamos na \olref[hal]{sec} por que exigir um
estado de parada designado não é uma restrição que impacta o que as
máquinas de Turing podem calcular.  Como as fitas de nossas máquinas de
Turing são infinitas em apenas uma direção, há casos em que uma máquina de
Turing não consegue executar adequadamente uma instrução: se ela lê o
quadrado mais à esquerda e deveria se mover para a esquerda. De acordo com a
nossa definição, ela simplesmente permanece no lugar em vez de ``cair'', mas
poderíamos tê-la definido de modo que ela pare quando isso acontecer. Essa
definição também é equivalente: poderíamos simular o comportamento de uma
máquina de Turing que para quando tenta se mover para a esquerda do
quadrado~$0$ excluindo cada transição $\delta(q,\TMendtape) =
\tuple{q',\sigma,\TMleft}$---então, em vez de tentar mover para a esquerda
em~$\TMendtape$, a máquina para.\footnote{Isso não funciona \emph{muito bem},
pois nada nos impede de escrever e ler $\TMendtape$
em quadrados que não sejam o quadrado~$0$ (veja \olref[una]{ex:mover}).
Podemos contornar isso adicionando um segundo símbolo~$\TMendtape'$ para
usar em seu lugar para tal propósito.}

Há também diferentes maneiras de representar números (e, portanto, a
função de entrada-saída calculada por uma máquina de Turing): usamos
representação unária, mas você também pode usar representação binária. Isso
requer dois símbolos além de $\TMblank$ e~$\TMendtape$.

Agora, aqui está um fato interessante: nenhuma dessas variações importa para
quais funções são Turing-computáveis. \emph{Se uma função é
Turing-computável segundo uma definição, ela é Turing-computável
segundo todas elas.}

Não entraremos nos detalhes de como verificar isso. Eis apenas um
exemplo: não ganhamos nenhum poder computacional adicional ao permitir uma
fita que é infinita em ambas as direções, ou múltiplas fitas. A razão é,
aproximadamente, que uma máquina de Turing com uma única fita infinita
unidirecional pode simular múltiplas fitas ou fitas infinitas bidirecionais.
Por exemplo, usando estados e instruções adicionais, podemos ``traduzir'' um
programa para uma máquina com múltiplas fitas ou fita infinita bidirecional
em um programa para uma máquina com uma única fita infinita unidirecional.
A máquina traduzida pode usar os quadrados pares para os quadrados da
fita~$1$ (ou os quadrados ``positivos'' de uma fita infinita bidirecional) e
os quadrados ímpares para os quadrados da fita~$2$ (ou os quadrados
``negativos'').

\end{document}
